% \iffalse meta-comment
%
%
% Copyright (C) 2017 by Mark W. Sprague
%
% This file may be distributed and/or modified under the conditions of
% the LaTeX Project Public License, either version 1.3 of this license
% or (at your option) any later version. The latest version of this
% license is in:
%
% http://www.latex-project.org/lppl.txt
%
% and version 1.3 or later is part of all distributions of LaTeX version
% 2005/12/01 or later.
%
% \fi
% \iffalse
%<*driver>
\ProvidesFile{ECULetter.dtx}
%</driver>
%<class>\NeedsTeXFormat{LaTeX2e}[2005/12/01]
%<class>\ProvidesClass{ECULetter}
%<*class>    
[2017/09/12 v2.0 ECU letterhead with letter and memo formatting]
%</class>
%
%<*driver>
\documentclass{ltxdoc}
\usepackage{hypdoc}
\EnableCrossrefs
\CodelineIndex
\RecordChanges
\begin{document}
  \DocInput{ECULetter.dtx}
\end{document}
%</driver>
% \fi
%
% \CheckSum{526}
%
% \CharacterTable
%  {Upper-case    \A\B\C\D\E\F\G\H\I\J\K\L\M\N\O\P\Q\R\S\T\U\V\W\X\Y\Z
%   Lower-case    \a\b\c\d\e\f\g\h\i\j\k\l\m\n\o\p\q\r\s\t\u\v\w\x\y\z
%   Digits        \0\1\2\3\4\5\6\7\8\9
%   Exclamation   \!     Double quote  \"     Hash (number) \#
%   Dollar        \$     Percent       \%     Ampersand     \&
%   Acute accent  \'     Left paren    \(     Right paren   \)
%   Asterisk      \*     Plus          \+     Comma         \,
%   Minus         \-     Point         \.     Solidus       \/
%   Colon         \:     Semicolon     \;     Less than     \<
%   Equals        \=     Greater than  \>     Question mark \?
%   Commercial at \@     Left bracket  \[     Backslash     \\
%   Right bracket \]     Circumflex    \^     Underscore    \_
%   Grave accent  \`     Left brace    \{     Vertical bar  \|
%   Right brace   \}     Tilde         \~}
%
%
% \newcommand{\hyperit}[1]{\textit{\hyperpage{#1}}}
%
% \changes{v1.0}{2017/03/21}{Initial version}
% \changes{v2.0}{2017/09/12}{ECU 2017 letterhead format}
%
% \GetFileInfo{ECULetter.dtx}
%
% \DoNotIndex{\newcommand,\newenvironment,\newif,\if,\else,\fi}
% \DoNotIndex{\setbox,\hbox,unskip,\ifdim,\wd}
% \DoNotIndex{\par,\vspace,\parskip}
% \DoNotIndex{\DeclareRobustCommand,\item,\unskip,\\}
%
% \title{The \textsf{ECULetter} class\thanks{This document
%   corresponds to \textsf{ECULetter}~\fileversion, dated \filedate.}}
% \author{Mark W. Sprague \\ \texttt{spraguem@ecu.edu}}
%
% \maketitle
%
% \begin{abstract}
% The \textsf{ECULetter} class provides a personalized ECU letterhead in
% a LaTeX environment. It also includes commands for elements of letters
% and memoranda written on the letterhead.
% \end{abstract}
%
% \tableofcontents
%
% \newpage
%
% \section{Introduction}
%
% The \textsf{ECULetter} class formats the document on official-looking
% ECU letterhead.  It also provides macros to format letters and memos
% on the letterhead.  User settings can be included in a file
% \index{ECULetterDefs.cfg file|hyperit} ECULetterDefs.cfg in the
% latex path.
%
% \section{Usage}
%
% \subsection{Class Declaration and Options}
%
% \DescribeMacro{\documentclass}
% The \index{class declaration|hyperit} class declaration statement is
% \begin{quote}
% |\documentclass|\oarg{options}|{ECULetter}|
% \end{quote}
% The following \index{class options|hyperit} options are valid:
% \begin{description}
% \item[|bw|] \index{bw|hyperit} Use black and white graphics for the
%             letterhead.
% \item[|color|] \index{color|hyperit} Use color graphics for the
%             letterhead (set by default).
% \item[|memo|] \index{memo|hyperit} Write the memo header to the
%             first page.
% \item[|noletterhead|] \index{noletterhead|hyperit} Do not write the
%             letterhead to the first page.
% \end{description}
%
% All other options are passed to the \textsf{article} class
% (\emph{n.b.}, not the letter class), upon which this class is based. 
% The \textsf{article} class options |10pt|, |12pt|, and |twocolumn| are
% not compatible with \textsf{ECULetter} and are silently ignored.
%
% \subsection{Required Packages}
%
% \changes{v2.0}{2017/09/12}{Package wrapfig no longer used}
% \changes{v2.0}{2017/09/12}{Package color required}
% \index{packages, required|hyperit} \index{geometry package|hyperit}
% The \textsf{ECULetter} class loads the packages |geometry|,
% |graphicx|, |color|, |fancyhdr|, |lastpage|, and |ragged2e|. 
% Commands from each of these packages may be used in the document.  Use
% the |\geometry| command from the |geometry| package to adjust the
% document layout and paper size (which defaults to |letterpaper|).
%
% \subsection{Required Graphics Files}
%
%  \changes{v2.0}{2017/09/12}{2017 ECU logo files}
% The \textsf{ECULetter} class requires three graphics files to be
% placed within the latex (|graphicx|) path: \index{ECULogo2017
% file|hyperit} |ECULogo2017| (color ECU logo, 2017 version), 
% \index{ECULogo2017BW file|hyperit} |ECULogo2017BW| (black and white ECU
% logo, 2017 edition), and
% \index{signature file|hyperit} |signature| (signature image).  These
% files may be in any format accessible to the
% |graphicx| package (with the appropriate extensions) and the files in
% the latex path may be symlinks.
%
%\subsection{Letterhead Content}
%
% The letterhead content is set with commands that begin with
% |\lh|\dots.  Omitting any of these definitions or defining it as empty
% (|{}|) or whitespace-only (|{  }|) results in that item being omitted
% from the letterhead.  The letterhead content is usually defined in the
% configuration file \index{ECULetterDefs.cfg file|hyperit}
% ECULetterDefs.cfg or in the document preamble.
%
% \DescribeMacro{\lhname}
% The author's name on the letterhead is set with
% \begin{quote}
%     |\lhname| \marg{Author Name}
% \end{quote}
% where $\langle$\emph{Author Name}$\rangle$ is the name as it appears
% at the top of the letterhead.  Remember to use a backslash (|\|) after
% periods following initials to prevent end-of-sentence spacing (for
% example, |John H.\ Doe|).
%
% \DescribeMacro{\lhdegree}
% The author's degree on the letterhead is set with
% \begin{quote}
%     |\lhdegree| \marg{Deg}
% \end{quote}
% where |Deg| is the degree abbreviation as it appears after the
% author's name on the letterhead.  The degree is only written to the
% letterhead if |\lhname| is defined.
%
% \DescribeMacro{\lhtitle}
% The author's title on the letterhead is set with
% \begin{quote}
%     |\lhtitle| \marg{Title}
% \end{quote}
% where $\langle$\emph{Title}$\rangle$ is the title (\emph{e.g.}
% Associate Professor) on the letterhead.  The title is only written to
% the letterhead if |\lhname| is defined.
%
% \DescribeMacro{\lhECUaddress}
% The department address on the letterhead is set with
% \begin{quote}
%     |\lhECUaddress| \marg{address}
% \end{quote}
% where $\langle$\emph{address}$\rangle$ is the campus portion of the
% address on the letterhead.  The $\langle$\emph{address}$\rangle$ can
% contain formatting and line breaks. The
% $\langle$\emph{address}$\rangle$ field can be used to format a
% departmental letterhead (with |\lhname| left undefined).
%
% \DescribeMacro{\lhECUPostal}
% The macro |\ECUPostal| gives the portion of the postal address after
% the departmental address.  This macro is defined as 
% \begin{quote}
%     East Carolina University \textbar\ Greenville, NC 27858-4353
% \end{quote}
% and can be redefined with %
% |\renewcommand{\lhECUPostal}|\marg{New Address}.
%
% \DescribeMacro{\lhphone}
% The telephone number on the letterhead is set with
% \begin{quote}
%     |\lhphone| \marg{xxx-xxx-xxxx}
% \end{quote}
% where $\langle$\emph{xxx-xxx-xxxx}$\rangle$ is the telephone number.
% Note that it can be formatted in as desired, \emph{e.g.} (123)
% 456-7890 or 123.456.7890.
%
% \DescribeMacro{\lhfax}
% The fax number on the letterhead is set with
% \begin{quote}
%     |\lhfax| \marg{xxx-xxx-xxxx}
% \end{quote}
% where $\langle$\emph{xxx-xxx-xxxx}$\rangle$ is the fax number. Note
% that it can be formatted in as desired, \emph{e.g.} (123)~456-7890 or
% 123.456.7890.
%
% \DescribeMacro{\lhemail}
% The E-mail address on the letterhead is set with
% \begin{quote}
%     |\lhemail| \marg{pirateid@ecu.edu}
% \end{quote}
% where $\langle$\emph{pirateid@ecu.edu}$\rangle$ is the E-mail address.
%
% \DescribeMacro{\lhwww}
% The web page URL on the letterhead is set with
% \begin{quote}
%     |\lhwww| \marg{ecuurl.ecu.edu}
% \end{quote}
% where $\langle$\emph{ecuurl.ecu.edu}$\rangle$ is the URL as should
% appear on the letterhead.
%
% \subsection{Memo Header Content}
%
% If the document is a memo, these commands, usually included in the
% document preamble, format the memo header.
%
% \DescribeMacro{\memotitle}
% The macro |\memotitle| gives the label MEMORANDUM at the top of the
% memo header and can be redefined with
% |\renewcommand{\memotitle}|\marg{newtitle}.
%
% The TO, FROM, DATE, and SUBJECT listings in the header are formatted
% using a |tabular| environment with the labels on the left and the
% appropriate fields on the right.
%
% \DescribeMacro{\memotolabel}
% The macro |\memotolabel| gives the label TO: in the memo header and
% can be redefined with
% |\renewcommand{\memotolabel}|\marg{newlabel}.
%
% \DescribeMacro{\memoto}
% The recipient of the memo is defined with
% \begin{quote}
%     |\memoto| \marg{therecipient}
% \end{quote}
% where $\langle$\emph{therecipient}$\rangle$ is the name listed in the
% TO-field of the memo header (TO: therecipient). Note that a second
% line (for example the recipient's title) can be included by specifying 
% \begin{quote}
%     |\memoto{|$\langle$\emph{therecipient}$\rangle$ |\\ &| %
%          $\langle$\emph{line 2}$\rangle$|}|
% \end{quote}
%
% \DescribeMacro{\memofromlabel}
% The macro |\memofromlabel| gives the label FROM: in the memo header
% and can be redefined with
% |\renewcommand{\memofromlabel}|\marg{newlabel}.
%
% \DescribeMacro{\memofrom}
% The sender of the memo is defined with
% \begin{quote}
%     |\memofrom| \marg{thesender}
% \end{quote}
% where $\langle$\emph{thesender}$\rangle$ is the name listed in the
% FROM-field of the memo header (FROM: thesender). Note that a second
% line (for example the sender's title) can be included by specifying 
% \begin{quote}
%     |\memofrom{|$\langle$\emph{thesender}$\rangle$ |\\ &| %
%          $\langle$\emph{line 2}$\rangle$|}|
% \end{quote}
% If |\memofrom| is not specified in the document, it defaults to the
% author's name on the letterhead as specified by |\lhname|.
%
% \DescribeMacro{\memodatelabel}
% The macro |\memodatelabel| gives the label DATE: in the memo header
% and can be redefined with
% |\renewcommand{\memodatelabel}|\marg{datelabel}.
%
% \DescribeMacro{\memodate}
% The date of the memo is defined with
% \begin{quote}
%     |\memodate| \marg{date}
% \end{quote}
% where $\langle$\emph{date}$\rangle$ is date to be used in in the DATE-field of the memo
% header (DATE: date). If |\memodate| is not supplied, it is replaced by
% |\today|.
%
% \DescribeMacro{\memosubjlabel}
% The macro |\memosubjlabel| gives the label SUBJECT: in the memo header
% and can be redefined with
% |\renewcommand{\memosubjlabel}|\marg{subjectlabel}.
%
% \DescribeMacro{\memosubject}
% The subject of the memo is defined with
% \begin{quote}
%     |\memosubject| \marg{subject text}
% \end{quote}
% where $\langle$\emph{subject text}$\rangle$ is used in in the
% SUBJECT-field of the memo header (SUBJECT: subject text). Normally the
% $\langle$\emph{subject text}$\rangle$ would fit in a one-line,
% regular-width table file.  A second line for the subject text can be
% included by specifying 
% \begin{quote}
%     |\memosubject{|$\langle$\emph{subject text}$\rangle$ %
%         |\\ &| $\langle$\emph{line 2}$\rangle$|}|
% \end{quote}
% Another possibility to extend the subject text is to place it inside a
% |minipage| environment.
%
% \subsection{Other Formatting}
% These preamble commands format other parts of the letter.
%
%  \changes{v2.0}{2017/09/12}{UNC statement removed from letterhead}
%
% \DescribeMacro{\ptwohead}
% An optional header included after the first page is defined with
% \begin{quote}
%     |\ptwohead| \marg{header text}
% \end{quote}
% If |\ptwohead| is not supplied, it defaults to the author's name on
% the letterhead as specified by |\lhname|.
%
% \subsection{Letter commands in the main document}
% These commands are used in the body of the document to format the
% letter.
%
% \DescribeMacro{\letterhead}
% The letterhead is written to the document using the content defined in
% the preamble with |\letterhead|. The |\letterhead| command should be
% the first thing written on the page.  This is done by default, unless
% the |noletterhead| class option is given. The |\letterhead| command
% allows the letterhead to be written to other pages (\emph{e.g.}, when
% the file contains multiple letters).
%
% \DescribeMacro{\memohead}
% The memo header is written to the document using the content defined
% in the preamble with |\memohead|. The memo header is written to the
% first page when the |memo| class option is given.  The |\memohead|
% command allows the memo header to be written to other pages.
%
% \DescribeMacro{\iaddress}
% The inside address of the letter is included with
% \begin{quote}
%     |\iaddress| \marg{address \textbackslash{}\textbackslash{} of \textbackslash{}\textbackslash{} recipient}
% \end{quote}
% This is usually the first command given in the body of the letter (but
% not for a memo) after the |\letterhrad| command.  The inside address
% is usually followed with the date, either using |\today| or the
% appropriate date text.
%
% \DescribeMacro{\greeting}
% The greeting in the letter is included with
% \begin{quote}
%     |\greeting| \marg{greeting text}
% \end{quote}
% For example, |\greeting{|Dear Jane:|}| would format the greeting
% text ``Dear Jane:'' on the left margin of the letter.  The |\greeting|
% command usually follows the date.
%
% \DescribeMacro{\closing}
% \DescribeMacro{\closingsig}
% \DescribeMacro{\closingname}
% The closing of the letter is given by |\closing| \marg{closing text}
% for a closing statement with name below and space for a handwritten
% signature.  The command |\closingsig| \oarg{sigwidth} \marg{closing
% text} inserts the graphic in the file \index{signature file|hyperit}
% |signature| into the signature space between the closing text and the
% closing name.  The optional argument
% $\langle$\emph{siggwidth}$\rangle$ sets the width of the signature
% graphic with default value |1.75in|.  The file 
% \index{signature file|hyperit} |signature| may be in any format
% accessible to the |graphicx| package.  The default name used below the
% closing is the letterhead name specified by |\lhname|, but it can be
% set with
% |\closingname| \marg{closing name}.  The $\langle$\emph{closing
% name}$\rangle$ may include formatting commands and line breaks (for
% titles, \emph{etc.}).
%
% \DescribeEnv{afterclose}
% \DescribeMacro{\ps}
% \DescribeMacro{\encl}
% \DescribeMacro{\cc}
% \DescribeMacro{\CC}
% \DescribeMacro{\pc}
% \DescribeMacro{\PC}
% The |afterclose| environment allows the letter postscript (P.S.),
% enclosure, and carbon copy statements .  The command |\ps| will format
% a postscript statement (P.S.) and format the text following it to have
% a hanging indent.  The command |\encl| will format an enclosure
% statement with the list of enclosed documents following the command. 
% A short enclosure list can be separated by commas, or a longer list
% can be separated by linebreaks.  The command |\cc| will format the
% carbon copy statement ``cc:'' and format the text following it to have
% a hanging indent.  Alternate copy statements |\CC| (for ``CC:''),
% |\pc| (for ``pc:''), and |\PC| (for ``PC:'') are also defined.  The
% |afterclose| environment is based on the latex |description|
% environment, so additional items can easily be provided with 
% |\item[|$\langle$\emph{label:}$\rangle$|]|.
%
% \subsection{User Definitions}
%
% The class reads user definitions from the file
% \index{ECULetterDefs.cfg file|hyperit} |ECULetterDefs.cfg| placed
% within the latex path.  Definitions in this file may be overridden on
% an \emph{ad-hoc} basis with commands in the document preamble. 
% Examples of useful user definitions include letterhead content
% definition commands so frequently-used values to not have to be
% specified in each document.
%
% \section{Sample Files}
%
% This section includes sample files for use with the \textsf{ECULetter}
% class.  The first sample file is a configuration file for storing
% default user settings.  The second sample file is a document template
% for a letter.  The third sample file is a document template for a memo.
%
% \subsection{Sample User Configuration File}
%
% \index{ECULetterDefs.cfg file|hyperit} 
% \begin{verbatim}
% % This is a sample user configuration file for use with the ECULetter
% % class.  Save this file with name ECULetterDefs.cfg, and save it
% % somewhere in the latex path.  The ECULetter class reads this file and
% % executes the commands within it at the end.  Add commands to define
% % letterhead parameters for the user.  Omit a command to leave the
% % parameter undefined, leaving it off of the letterhead.
% %
% % First define the name on the letterhead.
% \lhname{John Q.\ Doe}
% %
% % Define the degree: Ph.D., M.D., Ed.D., etc.
% \lhdegree{Ph.D.}
% %
% % Define the title: Professor, Head Honcho, etc. 
% \lhtitle{Associate Professor} 
% %
% % Define the departmental/campus address.
% \lhECUaddress{Department of Physics \textbar\ Mail Stop 563} 
% %
% % Define the telephone number. You can use different formats for the
% % number if desired. 
% \lhphone{252-867-5309}
% %
% % Define the fax number. Again, use the desired telephone number format.
% \lhfax{252-328-6314} 
% %
% % Define the E-mail address. 
% \lhemail{piratepd@ecu.edu}
% %
% % Define the web page.
% \lhwww{myweb.ecu.edu/piratepd}
% \end{verbatim}
%
% \subsection{Sample Document Template for Letter}
%
% \begin{verbatim}
% % This is a sample document template for a memo using the ECULetter class.
% \documentclass[ ]{ECULetter}
%
% \begin{document}
% \today
%
% % Insert inside address in brackets below.  Include linebreaks (\\).
% \iaddress{Address\\ of\\ Recipient} 
%
% % Insert the complete greeting with name and punctuation in brackets
% % below.
% \greeting{Dear Recipient:} 
%
% % Begin body of the letter below.
%
%
% % The following creates the closing. Include the closing text in the
% % brackets. Use \closing to sign by hand and \closingsig to use the
% % scanned signature in the file signature.pdf (stored in the latex path).
%
% \closing{Sincerely}
%
% % \closingsig{Sincerely}
%
% % Uncomment commands below to include afterclose items.
% % \begin{afterclose}
% %
% % \ps 
% %
% % \cc 
% %
% % \encl 
% %
% % \end{afterclose}
%
% \end{document}
% \end{verbatim}
%
% \subsection{Sample Document Template for Memo}
%
% \begin{verbatim}
% % This is a sample document template for a memo using the ECULetter class.
% \documentclass[memo]{ECULetter}
%
% % Define memo parameters.
% \memoto{Patricia T.\ Pirate} 
% % Include a two-line version with the following:
% % \memoto{Patricia T.\ Pirate, Chair\\ & Tailgating Committee}
%
% % The following defines the FROM field. Omit it to use the letterhead
% % name instead.
% \memofrom{P.\ D.\ the Pirate}
% % Include a two-line version with the following command.
% % \memofrom{P.\ D.\ the Pirate, Chair\\ & Trebuchet Committee}
%
% \memosubject{Saturday Activities}
%
% % The following defines the DATE field. Use any date format desired.
% % Omit the command to use \today.
% \memodate{October 31, 2017}
%
% \begin{document}
% %Begin body of the memo below.
%
%
% % Uncomment commands below to include afterclose items.
% % \begin{afterclose}
%
% % \ps
%
% % \cc 
%
% % \encl
%
% % \end{afterclose}
%
% \end{document}
% \end{verbatim}
%
% \section{Implementation}
%
% \StopEventually{^^A
%    \PrintChanges
%    \PrintIndex
% }
%
% Define boolean parameters to implement class options.
%    \begin{macrocode}
\newif\if@bw
\@bwfalse
%    \end{macrocode}
%
%    \begin{macrocode}
\newif\if@printletterhead
\@printletterheadtrue
%    \end{macrocode}
%
%    \begin{macrocode}
\newif\if@printmemohead
\@printmemoheadfalse
%    \end{macrocode}
%
%    \begin{macrocode}
\newif\if@sizegiven
\@sizegivenfalse
%    \end{macrocode}
% Declare options
%    \begin{macrocode}
\DeclareOption{bw}{\@bwtrue}
\DeclareOption{color}{\@bwfalse}
\DeclareOption{noletterhead}{\@printletterheadfalse}
\DeclareOption{memo}{\@printmemoheadtrue}
%    \end{macrocode}
%
%   \changes{v2.0}{2017/09/12}{Options 10pt and 12pt are now available}
%    The \textsf{ECULetter} class is based on the \textsf{article} class,
%    not the \textsf{letter} class.
%    Pass font size options to \textsf{article} class.
%    \begin{macrocode}
\DeclareOption{10pt}{
  \@sizegiventrue
  \PassOptionsToClass{10pt}{article}
}
\DeclareOption{11pt}{
  \@sizegiventrue
  \PassOptionsToClass{11pt}{article}
}
\DeclareOption{12pt}{
  \@sizegiventrue
  \PassOptionsToClass{12pt}{article}
}
%    \end{macrocode}
%
% Declare confliction \textsf{article} class options as |\OptionNotUsed|.
%    \begin{macrocode}
\DeclareOption{twocolumn}{\OptionNotUsed}
%    \end{macrocode}
%
% Pass all undefined options to the \textsf{article} class.
%    \begin{macrocode}
\DeclareOption*{%
  \PassOptionsToClass{\CurrentOption}{article}%
}
%    \end{macrocode}
%
% Process options.
%    \begin{macrocode}
\ProcessOptions\relax
%    \end{macrocode}
%
% If the font size option is not given, use 11pt as default instead of the 10pt default of
% the \textsf{article} class.
%    \begin{macrocode}
\if@sizegiven\else
  \PassOptionsToClass{11pt}{article}
\fi
%    \end{macrocode}
%
% Load the \textsf{article} class.
%    \begin{macrocode}
\LoadClass{article}
%    \end{macrocode}
%
% Load required packages.
%    \begin{macrocode}
\RequirePackage[letterpaper, hmargin={1in}, vmargin=1in]{geometry}
\RequirePackage{graphicx}
%    \end{macrocode}
%    \changes{v2.0}{2017/09/12}{Package color required}
%    \begin{macrocode}
\RequirePackage{color}
\RequirePackage{fancyhdr}
\RequirePackage{lastpage}
%    \end{macrocode}
%    \changes{v2.0}{2017/09/12}{Package wrapfig no longer used}
%    \begin{macrocode}
\RequirePackage{ragged2e}
%    \end{macrocode}
%
% Define ECU colors.
%    \changes{v2.0}{2017/09/12}{ECU colors for letterhead header, headrule, and footer}
%    \begin{macrocode}
\definecolor{ecupurple}{RGB}{80, 45, 127}
\definecolor{ecugold}{RGB}{255, 200, 0}
%    \end{macrocode}
% No headrule or footrule.
%    \begin{macrocode}
\renewcommand{\headrulewidth}{0pt}
\renewcommand{\footrulewidth}{0pt}
%    \end{macrocode}
%
% Take care of letterhead definitions.  Most commands have a boolean
% |\if@commanddefined| and a value holder |\@command|.
%
%    \begin{macrocode}
\newif\if@lhnamedefined
\@lhnamedefinedfalse
%    \end{macrocode}
%
%    \begin{macrocode}
\DeclareRobustCommand*{\@lhname}{}
\DeclareRobustCommand*{\lhname}[1]{
  \DeclareRobustCommand*{\@lhname}{#1}%
%    \end{macrocode}
% The following command |\setbox|\dots\ strips whitespace and declares
% the command undefined if nothing is left. It will cause an error if
% there are any |\\| in the value (use |\linebreak| instead).
%    \begin{macrocode}
  \setbox0=\hbox{\@lhname\unskip}\ifdim\wd0=0pt
    \@lhnamedefinedfalse
  \else
    \@lhnamedefinedtrue
  \fi
}
%    \end{macrocode}
%
%    \begin{macrocode}
\newif\if@lhdegreedefined
\@lhdegreedefinedfalse
%    \end{macrocode}
%    \begin{macrocode}
\DeclareRobustCommand*{\@lhdegree}{}
\DeclareRobustCommand*{\lhdegree}[1]{
  \DeclareRobustCommand*{\@lhdegree}{#1}
  \setbox0=\hbox{\@lhdegree\unskip}\ifdim\wd0=0pt
    \@lhdegreedefinedfalse
  \else
    \@lhdegreedefinedtrue
  \fi
}
%    \end{macrocode}
%
%    \begin{macrocode}
\newif\if@lhtitledefined
\@lhtitledefinedfalse
%    \end{macrocode}
%    \begin{macrocode}
\DeclareRobustCommand*{\@lhtitle}{}
\DeclareRobustCommand*{\lhtitle}[1]{
  \DeclareRobustCommand*{\@lhtitle}{#1}
  \setbox0=\hbox{\@lhtitle\unskip}\ifdim\wd0=0pt
    \@lhtitledefinedfalse
  \else
    \@lhtitledefinedtrue
  \fi
}
%    \end{macrocode}
%
%    \begin{macrocode}
\newif\if@lhECUaddressdefined
\@lhECUaddressdefinedfalse
%    \end{macrocode}
%    \begin{macrocode}
\DeclareRobustCommand{\@lhECUaddress}{}
\DeclareRobustCommand{\lhECUaddress}[1]{
  \DeclareRobustCommand*{\@lhECUaddress}{#1}
  \@lhECUaddressdefinedtrue
}
%    \end{macrocode}
%
%    \begin{macrocode}
\DeclareRobustCommand{\lhECUPostal}{
    East Carolina University \textbar\ 
    Greenville, NC 27858-4353
}
%    \end{macrocode}
%
%    \begin{macrocode}
\newif\if@lhphonedefined
\@lhphonedefinedfalse
%    \end{macrocode}
%    \begin{macrocode}
\DeclareRobustCommand*{\@lhphone}{}
\DeclareRobustCommand*{\lhphone}[1]{
  \DeclareRobustCommand*{\@lhphone}{#1}
  \setbox0=\hbox{\@lhphone\unskip}\ifdim\wd0=0pt
    \@lhphonedefinedfalse
  \else
    \@lhphonedefinedtrue
  \fi
}
%    \end{macrocode}
%
%    \begin{macrocode}
\newif\if@lhfaxdefined
\@lhfaxdefinedfalse
%    \end{macrocode}
%    \begin{macrocode}
\DeclareRobustCommand*{\@lhfax}{}
\DeclareRobustCommand*{\lhfax}[1]{
  \DeclareRobustCommand*{\@lhfax}{#1}
  \setbox0=\hbox{\@lhfax\unskip}\ifdim\wd0=0pt
    \@lhfaxdefinedfalse
  \else
    \@lhfaxdefinedtrue
  \fi
}
%    \end{macrocode}
%
%    \begin{macrocode}
\newif\if@lhemaildefined
\@lhemaildefinedfalse
%    \end{macrocode}
%    \begin{macrocode}
\DeclareRobustCommand*{\@lhemail}{}
\DeclareRobustCommand*{\lhemail}[1]{
  \DeclareRobustCommand*{\@lhemail}{#1}
  \setbox0=\hbox{\@lhemail\unskip}\ifdim\wd0=0pt
    \@lhemaildefinedfalse
  \else
    \@lhemaildefinedtrue
  \fi
}
%    \end{macrocode}
%
%    \begin{macrocode}
\newif\if@lhwwwdefined
\@lhwwwdefinedfalse
%    \end{macrocode}
%    \begin{macrocode}
\DeclareRobustCommand*{\@lhwww}{}
\DeclareRobustCommand*{\lhwww}[1]{
  \DeclareRobustCommand*{\@lhwww}{#1}
  \setbox0=\hbox{\@lhwww\unskip}\ifdim\wd0=0pt
    \@lhwwwdefinedfalse
  \else
    \@lhwwwdefinedtrue
  \fi
}
%    \end{macrocode}
%
% Declare memo commands
%
%    \begin{macrocode}
\DeclareRobustCommand*{\memotitle}{MEMORANDUM}
%    \end{macrocode}
%
%    \begin{macrocode}
\DeclareRobustCommand*{\memotolabel}{TO:}
%    \end{macrocode}
%
%    \begin{macrocode}
\DeclareRobustCommand*{\@memoto}{} 
\DeclareRobustCommand*{\memoto}[1]{\DeclareRobustCommand*{\@memoto}{#1}} 
%    \end{macrocode}
%
%    \begin{macrocode}
\newif\if@memofromdefined
\@memofromdefinedfalse
%    \end{macrocode}
%
%    \begin{macrocode}
\DeclareRobustCommand*{\memofromlabel}{FROM:}
%    \end{macrocode}
%
%    \begin{macrocode}
\DeclareRobustCommand*{\@memofrom}{} 
\DeclareRobustCommand*{\memofrom}[1]{
  \DeclareRobustCommand*{\@memofrom}{#1} 
  \@memofromdefinedtrue
}
%    \end{macrocode}
%
%    \begin{macrocode}
\DeclareRobustCommand*{\memosubjlabel}{SUBJECT:}
%    \end{macrocode}
%
%    \begin{macrocode}
\DeclareRobustCommand*{\@memosubject}{} 
\DeclareRobustCommand*{\memosubject}[1]{\DeclareRobustCommand*{\@memosubject}{#1}}
%    \end{macrocode}
%
%    \begin{macrocode}
\newif\if@memodatedefined
\@memodatedefinedfalse
%    \end{macrocode}
%
%    \begin{macrocode}
\DeclareRobustCommand*{\memodatelabel}{DATE:}
%    \end{macrocode}
%
%    \begin{macrocode}
\DeclareRobustCommand*{\@memodate}{} 
\DeclareRobustCommand*{\memodate}[1]{
  \DeclareRobustCommand*{\@memodate}{#1} 
  \@memodatedefinedtrue
}
%    \end{macrocode}
%
% Declare other formatting commands.
%
%    \begin{macrocode}
\newif\if@ptwoheaddefined
\@ptwoheaddefinedfalse
%    \end{macrocode}
%
%    \begin{macrocode}
\DeclareRobustCommand*{\@ptwohead}{}
\DeclareRobustCommand*{\ptwohead}[1]{
  \DeclareRobustCommand*{\@ptwohead}{#1}
  \@ptwoheaddefinedtrue
}
%    \end{macrocode}
%
%    \changes{v2.0}{2017/09/12}{UNC statement removed from letterhead}
%
%    \begin{macrocode}
\DeclareRobustCommand{\iaddress}[1]{
  \par
  #1
  \par
}
%    \end{macrocode}
%
%    \begin{macrocode}
\DeclareRobustCommand{\greeting}[1]{
  \par
  #1
  \par
}
%    \end{macrocode}
%
%    \begin{macrocode}
\newif\if@closingnamedefined
\@closingnamedefinedfalse
%    \end{macrocode}
%
%    \begin{macrocode}
\DeclareRobustCommand*{\@closingname}{}
\DeclareRobustCommand*{\closingname}[1]{
  \DeclareRobustCommand*{\@closingname}{#1}
  \@closingnamedefinedtrue
}
%    \end{macrocode}
%
%    \begin{macrocode}
\DeclareRobustCommand{\closing}[1]{
  \par
  \vspace{\parskip}\vspace{\parskip}#1,
  \par
  \vspace{\parskip}\vspace{\parskip}%
  \if@closingnamedefined
    \@closingname
  \else
    \if@lhnamedefined
      \@lhname
    \else
      ~
    \fi
  \fi
}
%    \end{macrocode}
%
%    \begin{macrocode}
\newlength{\sigwidth}
\setlength{\sigwidth}{1.75in}
%    \end{macrocode}
%
% \begin{macro}{\closingsig}
% This command will include a signature image in the file
% \index{signature file|hyperit} |signature|.
%    \begin{macrocode}
\DeclareRobustCommand{\closingsig}[2][\sigwidth]{
  \par
  \vspace{\parskip}\vspace{\parskip}#2,\\
  \includegraphics[width=#1]{signature}\\
  \if@closingnamedefined
    \@closingname
  \else
    \if@lhnamedefined
      \@lhname
    \else
      ~
    \fi
  \fi
}
%    \end{macrocode}
%\end{macro}
%
%\begin{macro}{\letterhead}
% This is the command to include the letterhead in the document.
%    \begin{macrocode}
\DeclareRobustCommand{\letterhead}{
  \thispagestyle{lhpage}
}
%    \end{macrocode}
% \end{macro}
%
% \begin{macro}{\memohead}
% The command memohead creates a heading for a memorandum.
%    \begin{macrocode}
\DeclareRobustCommand{\memohead}{
\begin{tabular}{l l}
  \multicolumn{2}{l}{\memotitle}\\ \\
  \memotolabel & \@memoto\\ \\
  \memofromlabel & 
  \if@memofromdefined
    \@memofrom\\ \\
  \else
    \if@lhnamedefined
      \@lhname\\ \\
    \else
      ~\\ \\
    \fi
  \fi
  \memodatelabel & 
  \if@memodatedefined
    \@memodate \\ \\
  \else
    \today \\ \\
  \fi
  \memosubjlabel & \@memosubject\\ \\
\end{tabular}
}
%    \end{macrocode}
% \end{macro}
%
% The following length definitions are necessary.
%    \begin{macrocode}
\headsep = 0.25in
\parskip = 2ex
\parindent = 0.0in
\footskip = 18.7pt
%    \end{macrocode}
%%
% \begin{environment}{afterclose}
% Define the |afterclose| environment.
%    \begin{macrocode}
\newenvironment{afterclose}{\begin{description}}{\end{description}}
%    \end{macrocode}
% \end{environment}
%
% Define the |afterclose| items.
%    \begin{macrocode}
\DeclareRobustCommand*{\ps}{
  \item[{\rm{P.S.}}] 
}
%    \end{macrocode}
%
%    \begin{macrocode}
\DeclareRobustCommand*{\encl}{
  \item[{\rm{encl:}}] 
}
%    \end{macrocode}
%
%    \begin{macrocode}
\DeclareRobustCommand*{\cc}{
  \item[{\rm{cc:}}] 
}
%    \end{macrocode}
%
%    \begin{macrocode}
\DeclareRobustCommand*{\CC}{
  \item[{\rm{CC:}}] 
}
%    \end{macrocode}
%
%    \begin{macrocode}
\DeclareRobustCommand*{\pc}{
  \item[{\rm{pc:}}] 
}
%    \end{macrocode}
%
%    \begin{macrocode}
\DeclareRobustCommand*{\PC}{
  \item[{\rm{PC:}}] 
}
%    \end{macrocode}
% Define the page style for the letterhead page.
%    \begin{macrocode}
\fancypagestyle{lhpage}{%
     \newgeometry{vmargin={2in, 1in}, hmargin=1in, headheight=62pt, foot=6pt}
%    \end{macrocode}
%    \changes{v2.0}{2017/09/12}{Headrule in letterhead}
%    \begin{macrocode}
     \renewcommand{\headrulewidth}{2pt}
     \renewcommand{\headrule}{\hbox to\headwidth{%
       \if@bw\else
         \color{ecugold}%
       \fi
       \leaders\hrule height \headrulewidth\hfill}}
%    \end{macrocode}
% The black and white  and color ECU logo files are in the |\includegraphics| commands below. 
% these files must be somewhere in the latex path. 
% (Test with the terminal command |$ kpsewhich filename.pdf|.)
%    \begin{macrocode}
  \fancyhf{}
  \fancyhead[L]{%
    \if@bw
      \includegraphics[scale=0.75, viewport=48 8 144 72]{ECULogo2017BW}
    \else
      \includegraphics[scale=0.75, viewport=48 8 144 72]{ECULogo2017}
    \fi
  }
%    \end{macrocode}
% Next the letterhead text. For a color letterhead, use the color ecupurple defined above.
%    \begin{macrocode}
  \fancyhead[R]{
    \if@bw\else
      \color{ecupurple}
    \fi
    \sf\raggedleft \fontsize{9}{10.8}\selectfont
     \if@lhnamedefined
       \textbf{\@lhname\if@lhdegreedefined{, \@lhdegree}\fi}
          \if@lhtitledefined
            \textbar\ \@lhtitle\\
          \else
            \\
          \fi
        \fi
      \if@lhECUaddressdefined
        \@lhECUaddress \textbar 
      \fi
      \lhECUPostal\\
      \if@lhphonedefined
        \@lhphone\ office
      \fi
      \if@lhfaxdefined
        \if@lhphonedefined
          \textbar\ 
        \fi
        \@lhfax\ fax
      \fi
      \if@lhemaildefined
        \textbar\ \@lhemail\ 
      \fi
      \if@lhwwwdefined
        \textbar\ \@lhwww
      \fi
  }
  \fancyfoot{}
  \fancyfoot[R]{
    \if@bw\else
      \color{ecugold}
    \fi
    \sf\raggedleft \fontsize{16}{19.2}\selectfont
    \vspace{-12pt}
    www.ecu.edu
  }
}
%    \end{macrocode}
% Set the page style for pages following the first page.
%    \begin{macrocode}
\pagestyle{fancy}
\newgeometry{vmargin=1in, hmargin=1in, headheight=16pt, includehead}
\renewcommand{\headrulewidth}{0pt}
\lhead{
  \if@ptwoheaddefined 
    \@ptwohead
  \else
    \@lhname
  \fi
}
\chead{}
\rhead{Page \thepage\ of \pageref{LastPage}}
\lfoot{}
\cfoot{}
\rfoot{}
%    \end{macrocode}
%
% Set the main text to be |\raggedright|
%    \begin{macrocode}
\raggedright
%    \end{macrocode}
%
% Read user defined commands from the file 
% \index{ECULetterDefs.cfg file|hyperit} |ECULetterDefs.cfg| if it exists.
%    \begin{macrocode}
\InputIfFileExists{ECULetterDefs.cfg}{%
  \ClassInfo{ECULetter}{Loading user configurations from ECULetterDefs.cfg}%
}%
{%
  \ClassInfo{ECULetter}{User configuration file found.}%
}
%    \end{macrocode}
%
% Include |\letterhead| and |\memohead| at the beginning of the document
% if appropriate.
%    \begin{macrocode}
\AtBeginDocument{
  \if@printletterhead
    \letterhead
  \fi
  \if@printmemohead
    \memohead
    \par
  \fi
}
%    \end{macrocode}
%
%
% \Finale
\endinput
